% !TEX program = pdflatex
\documentclass[11pt,a4paper]{article}

\usepackage[utf8]{inputenc}
\usepackage[T1]{fontenc}
\usepackage{lmodern}
\usepackage[margin=2.5cm]{geometry}
\usepackage{amsmath,amssymb,amsthm}
\usepackage{booktabs}
\usepackage{enumitem}
\usepackage{microtype}
\usepackage{hyperref}

\hypersetup{colorlinks=true,linkcolor=black,citecolor=black,urlcolor=blue!50!black}

\newtheorem{theorem}{Theorem}
\newtheorem{lemma}[theorem]{Lemma}
\newtheorem{corollary}[theorem]{Corollary}
\theoremstyle{definition}
\newtheorem{definition}[theorem]{Definition}
\theoremstyle{remark}
\newtheorem{remark}[theorem]{Remark}

\newcommand{\TT}{T}
\newcommand{\Pleft}{\mathcal{P}_{\mathrm{left}}}
\newcommand{\Pland}{\mathcal{P}_{\mathrm{landing}}}
\newcommand{\Pright}{\mathcal{P}_{\mathrm{right}}}

\title{A Lean~4 Formalization of a Bertrand--Syracuse Prime Trichotomy}
\author{Thomas Hofbauer\\
\small Bamberg / Akademie Olympia\\
\small\texttt{Kepler--Hurwitz}}
\date{28 July 2026}

\begin{document}
\maketitle

\begin{abstract}
We formalize a local interaction between prime distribution in Bertrand
intervals and the odd-step Syracuse map. For every odd integer $n>1$
satisfying $v_2(3n+1)=1$, the Syracuse successor
\[
\TT(n)=\frac{3n+1}{2}
\]
lies strictly inside the interval $(n,2n]$. It therefore determines a
canonical partition of the prime spectrum in this interval into primes
lying to the left of $\TT(n)$, the possible prime landing point $\TT(n)$
itself, and primes lying to its right. We prove that these classes are
pairwise disjoint and exhaust all primes in the Bertrand interval. The
landing class is a singleton exactly when $\TT(n)$ is prime and is empty
otherwise.

The result provides a formally verified bridge between a bounded
prime-number window and the local orbit geometry of the Syracuse
transformation. It is a finite classification theorem and makes no claim
of proving global termination of Collatz orbits. The complete development
has been checked in Lean~4 without unproved placeholders.
\end{abstract}

\section{Introduction}

Two classical objects meet in a finite window:
\[
(n,2n]\cap\mathbb{P}
\qquad\text{and}\qquad
\TT(n)=\frac{3n+1}{2^{v_2(3n+1)}}.
\]
We study neither the asymptotic distribution of primes nor the global
termination of Collatz orbits. Our result concerns a canonical finite
partition induced by a single odd Syracuse step.

Under the standing condition $v_2(3n+1)=1$, one has simply
$\TT(n)=(3n+1)/2$. For odd $n>1$ this successor lies in the interior of
the Bertrand window:
\[
n<\TT(n)<2n.
\]
Thus $\TT(n)$ is not a free threshold. It is the actual local orbit
successor of~$n$, and it cuts the Bertrand prime spectrum at an
\emph{orbit-determined partition point}
(equivalently: a distinguished dynamical landing point).
\begin{remark}
In the terminology of the broader EABC framework, this point may be
regarded as a local resonance point. That interpretive language is not
required for the theorems below and is not used as a primary technical
term.
\end{remark}

\begin{remark}[Detachability]
This note is intentionally independent of any Dumas-holography wording and
of any affine-transport principle. Those larger bridges remain separate
programme layers; the present theorem does not import them.
\end{remark}

\section{Arithmetic preliminaries}

\begin{definition}
Write $\TT(n):=(3n+1)/2$ whenever $v_2(3n+1)=1$.
\end{definition}

\begin{lemma}[Residue characterisation]
\label{lem:val-mod}
Let $n$ be odd. Then
\[
v_2(3n+1)=1 \iff n\equiv 3\pmod{4}.
\]
\end{lemma}
\begin{proof}[Sketch]
Since $n$ is odd, $n\equiv 1$ or $3\pmod{4}$. If $n\equiv 1\pmod{4}$, then
$3n+1\equiv 0\pmod{4}$, so $v_2(3n+1)\ge 2$. If $n\equiv 3\pmod{4}$, then
$3n+1=2(6k+5)$ with $6k+5$ odd, hence $v_2(3n+1)=1$.
\end{proof}
Thus the family under study is exactly one natural half of the odd
residue classes, not an artificial valuation slice.

\begin{definition}
The Bertrand prime spectrum of $n$ is
\[
\mathbb{P}\cap(n,2n]
=
\{p\in\mathbb{P}: n<p\le 2n\}.
\]
\end{definition}

\section{Interior landing lemma}

\begin{lemma}[Interior landing]
\label{lem:interior}
Let $n>1$ and $n\equiv 3\pmod{4}$. Then
\[
n<\TT(n)<2n.
\]
\end{lemma}
\begin{proof}
The hypothesis yields $2\TT(n)=3n+1$. Then
\[
\TT(n)-n=\frac{n+1}{2}>0,
\qquad
2n-\TT(n)=\frac{n-1}{2}>0.
\]
\end{proof}
This inequality is the hinge between dynamics and interval arithmetic. It
precedes any prime partition: the reader sees immediately that the primes
in the Bertrand window are cut at a dynamically determined point, that the
point is not a free threshold, that it is the local orbit successor of~$n$,
and that the statement remains fully local and finite.

\section{Prime-spectrum trichotomy}

\begin{definition}
\[
\begin{aligned}
\Pleft(n)
&=\{p\in\mathbb{P}: n<p<\TT(n)\},\\
\Pland(n)
&=\{p\in\mathbb{P}: p=\TT(n)\},\\
\Pright(n)
&=\{p\in\mathbb{P}: \TT(n)<p\le 2n\}.
\end{aligned}
\]
\end{definition}
In particular,
\[
\Pland(n)
=
\begin{cases}
\{\TT(n)\}, & \TT(n)\in\mathbb{P},\\
\varnothing, & \TT(n)\notin\mathbb{P}.
\end{cases}
\]
Hence $\#\Pland(n)\le 1$.

\begin{theorem}[Bertrand--Syracuse prime trichotomy]
\label{thm:main}
Let $n>1$ and $n\equiv 3\pmod{4}$. Define $\TT(n)=(3n+1)/2$.
Then $n<\TT(n)<2n$. Consequently the primes in the Bertrand interval admit
the disjoint decomposition
\[
\mathbb{P}\cap(n,2n]
=
\Pleft(n)
\,\dot\cup\,
\Pland(n)
\,\dot\cup\,
\Pright(n),
\]
where the three classes are as above.
Moreover $\Pland(n)=\{\TT(n)\}$ if and only if $\TT(n)$ is prime, and
$\Pland(n)=\varnothing$ otherwise.
\end{theorem}
\begin{proof}[Proof outline]
Lemma~\ref{lem:interior} places $\TT(n)$ in $(n,2n)$. For any prime
$p\in(n,2n]$, the trichotomy $p<\TT(n)$, $p=\TT(n)$, or $p>\TT(n)$ yields
membership in exactly one of $\Pleft(n)$, $\Pland(n)$, $\Pright(n)$.
Disjointness is immediate from the strict inequalities. The landing
characterisation is the definition of $\Pland(n)$.
\end{proof}

Two layers must be kept distinct:
\begin{enumerate}[leftmargin=*]
  \item \textbf{Order-theoretic trichotomy.} The partition follows from the
  linear order on~$\mathbb{N}$ once Lemma~\ref{lem:interior} places
  $\TT(n)$ inside $(n,2n]$. Bertrand's postulate is not used here.
  \item \textbf{Bertrand non-emptiness.} Bertrand's postulate supplies
  $\mathbb{P}\cap(n,2n]\neq\varnothing$ for $n>1$. The name
  ``Bertrand--Syracuse'' therefore combines an order partition with an
  independent existence theorem; Bertrand is not credited with proving the
  trichotomy itself.
\end{enumerate}

\section{Bertrand consequences}

\begin{corollary}
\label{cor:landing-empty}
If $\Pland(n)=\varnothing$, then
$\Pleft(n)\neq\varnothing$ or $\Pright(n)\neq\varnothing$.
\end{corollary}
\begin{proof}
By Bertrand, $\mathbb{P}\cap(n,2n]\neq\varnothing$. Theorem~\ref{thm:main}
exhausts this set by the three classes; with empty landing class, at least
one of the remaining classes is nonempty.
\end{proof}

\begin{corollary}
\label{cor:unique-landing}
If $\TT(n)$ is the unique prime in $(n,2n]$, then
$\Pleft(n)=\Pright(n)=\varnothing$ and $\Pland(n)=\{\TT(n)\}$.
\end{corollary}
\begin{proof}
Uniqueness forces $\mathbb{P}\cap(n,2n]=\{\TT(n)\}$, so $\TT(n)$ is prime
and no prime lies strictly left or right of it inside the window.
\end{proof}

\section{Scope and non-claims}

The formalization does not establish that the Syracuse map is globally
terminating. It does not derive monotone descent, bounded stopping times,
or a global covering of odd integers by descending cylinders. The
construction records how one local orbit step canonically organizes the
primes in a bounded interval.

In particular, a schema-level equivalence of the form
``totality of a formal predicate $\Leftrightarrow$ Collatz''---if present
elsewhere in a larger development---is a characterisation of global reach,
not a proof that the predicate holds. No such claim is made here.

Interpretive vocabulary (``resonance'', ``arithmetico-dynamical coupling'',
``local spectral structure'', EABC geometry) is optional commentary only.

\section{Formalization architecture}

The reproducible preprint package lives at
\texttt{docs/manuscripts/bertrand\_syracuse\_preprint\_package/}
(README, \texttt{MANIFEST.json}, \texttt{reproduce.sh}).

\begin{center}
\begin{tabular}{@{}ll@{}}
\toprule
Item & Value \\
\midrule
Lean & 4.31.0 \\
Mathlib & \texttt{v4.31.0} \\
Module &
\texttt{KeplerHurwitz.Collatz.BertrandSyracuseTrichotomy} \\
Build &
\texttt{lake build KeplerHurwitz.Collatz.BertrandSyracuseTrichotomy} \\
Package check &
\texttt{./docs/manuscripts/bertrand\_syracuse\_preprint\_package/reproduce.sh} \\
Sorry audit &
\texttt{rg -n -P '(sorry\textbar admit)'
  KeplerHurwitz/Collatz/BertrandSyracuseTrichotomy.lean} \\
Certificate &
\texttt{bertrandSyracuseTrichotomyCertificate} \\
Placeholders & \texttt{0 sorry} / \texttt{0 admit} in the module \\
Seal & PREPRINT-READY $\cdot$ FORMAL CORE SEALED \\
\bottomrule
\end{tabular}
\end{center}

Principal Lean names:
\texttt{padicValNat\_two\_three\_n\_succ\_eq\_one\_iff},
\texttt{syracuseT\_mem\_interior},
\texttt{bertrandPrimeSpectrum\_eq\_union},
\texttt{landing\_eq\_singleton\_iff\_prime},
\texttt{landing\_ncard\_le\_one},
\texttt{bertrandPrimeSpectrum\_nonempty},
\texttt{left\_or\_right\_of\_landing\_empty},
\texttt{unique\_prime\_is\_landing}.

\section*{Acknowledgements}
The author thanks the classical $3x+1$ tradition---Collatz, Terras,
Lagarias, and many others---for the language in which a finite local
statement can be stated without overclaiming a global solution.

\begin{thebibliography}{9}
\bibitem{Bertrand}
J.~Bertrand.
\newblock M\'emoire sur le nombre de valeurs que peut prendre une fonction
qui d\'epend de plusieurs variables.
\newblock \emph{J.\ \'Ec.\ Polytech.} \textbf{30} (1845), 123--140.

\bibitem{Erdos1932}
P.~Erd\H{o}s.
\newblock Beweis eines Satzes von Tschebyschef.
\newblock \emph{Acta Sci.\ Math.\ (Szeged)} \textbf{5} (1932), 194--198.

\bibitem{Terras1976}
R.~Terras.
\newblock A stopping time problem on the positive integers.
\newblock \emph{Acta Arith.} \textbf{30} (1976), 241--252.

\bibitem{Lagarias2010}
J.~C.~Lagarias (ed.).
\newblock \emph{The Ultimate Challenge: The $3x+1$ Problem}.
\newblock American Mathematical Society, 2010.

\bibitem{MathlibBertrand}
The mathlib Community.
\newblock Bertrand's postulate
(\texttt{Mathlib.NumberTheory.Bertrand}).
\end{thebibliography}

\end{document}
